\chapter{Y}
\label{chap:y}

\term{Yang-Mills Theory}
A gauge theory that generalizes electromagnetism to non-abelian Lie groups. It provides the mathematical framework for the strong and weak nuclear forces in the Standard Model of particle physics. The existence of a "mass gap" in Yang-Mills theory is one of the Millennium Prize Problems.

\term{Yang-Baxter Equation}
A consistency equation for a system of three particles undergoing pairwise scattering, $R_{12}R_{13}R_{23} = R_{23}R_{13}R_{12}$. It is central to the study of integrable systems, quantum groups, and knot invariants.

\term{y-axis}
The vertical coordinate axis in the Cartesian plane, usually measuring the second coordinate.

\term{y-coordinate}
The second coordinate of a point, measuring its signed distance from the x-axis.

\term{y-intercept}
The point where a line or curve crosses the y-axis, occurring where $x = 0$.

\term{Young Diagram}
A finite collection of boxes arranged in left-aligned rows, with the row lengths in non-increasing order. They provide a visual representation of integer partitions and index the irreducible representations of the symmetric group $S_n$.

\term{Young Tableau}
A filling of a Young diagram with positive integers. A tableau is "standard" if the entries are strictly increasing across each row and down each column. The number of standard Young tableaux of a given shape can be computed using the Hook Length Formula.

\term{Young's Inequality}
A result in analysis stating that for $a, b \geq 0$ and $1/p + 1/q = 1$ (with $p, q > 1$), $ab \leq \frac{a^p}{p} + \frac{b^q}{q}$. This is a generalization of the AM-GM inequality and is a key step in proving Hölder's inequality.

\term{Yoneda Lemma}
A fundamental result in category theory stating that for any object $A$ in a category $\mathcal{C}$, the natural transformations from the representable functor $\operatorname{Hom}(A, -)$ to any functor $F: \mathcal{C} \to \text{Set}$ are in bijection with the elements of $F(A)$. It implies that an object is completely determined by its relationships to other objects.

\term{Y-Delta Transform}
A method in network theory for simplifying electrical circuits by replacing a "star" (Y) configuration of three resistors with an equivalent "delta" ($\Delta$) configuration, and vice versa.
